% Figure: available work from a finite hot body cooling to ambient. As the body
% cools from its initial temperature toward the dead state, the maximum work
% recoverable per unit heat shed falls, because each later joule is shed at a
% lower temperature and so carries a smaller Carnot factor. The curve plots the
% running exergy still extractable against the body's current temperature; the
% area interpretation is the integral of (1 - T0/T) dQ.
\begin{figure}[htbp]
\centering
\begin{tikzpicture}
\begin{axis}[
  width=11cm, height=7cm,
  xlabel={body temperature $T$ during cooling ($\si{\kelvin}$)},
  ylabel={local Carnot factor $1 - T_0/T$},
  xmin=290, xmax=600, ymin=0, ymax=0.55,
  axis lines=left,
  tick label style={font=\scriptsize},
  label style={font=\small},
  legend style={font=\scriptsize, at={(0.03,0.97)}, anchor=north west, draw=enggray!40},
  clip=false,
]
% local Carnot factor for T0 = 290 K
\addplot[engblue, very thick, domain=290:600, samples=120] {1 - 290/x};
\addlegendentry{$1 - T_0/T$, $T_0 = 290\,\si{\kelvin}$}
% shade the region whose area (against dQ = mc dT) is the recoverable work
\addplot[draw=none, fill=engblue!12, domain=290:600, samples=80] {1 - 290/x} \closedcycle;
% mark the initial state
\draw[enggray, dashed] (axis cs:600,0) -- (axis cs:600,{1-290/600});
\node[engblue, font=\scriptsize, anchor=south west] at (axis cs:560,{1-290/560})
  {start at $T_i = 600\,\si{\kelvin}$};
\node[enggray, font=\scriptsize, anchor=south] at (axis cs:330,0.03)
  {dead state};
\end{axis}
\end{tikzpicture}
\caption[Available work from a finite hot body cooling to ambient]{The Carnot
factor $1 - T_0/T$ that weights each increment of heat shed by a hot body as it
cools from its initial temperature $T_i$ toward the dead state $T_0$. A finite
body is not a reservoir: as it cools, every later joule leaves at a lower
temperature and so converts to work at a smaller Carnot factor, until at the dead
state the factor reaches zero and no further work can be drawn. The maximum work
recoverable is the integral of this factor against the heat shed, $\int (1 -
T_0/T)\,mc\,dT$ from $T_0$ to $T_i$, the shaded area read against the body's heat
capacity. The curve makes plain why a finite hot stream yields less work per joule
than a true high-temperature reservoir: the reservoir holds its temperature while
the finite body slides down the curve, and the available work is the average of
the factor over the descent, not its value at the start.}
\label{fig:vol04ch04:finite-reservoir-exergy}
\end{figure}
